\scanpage{51}
% loc: heading above the rule
Remarque sur la caractérisation des classes de coh. D.R. \deleted{\unsure{qu}}\inserted{\unsure{essentiellement}} \uline{algébriques.} par bigraduation ou par réalité \unsure{....}
\par\noindent\rule{\linewidth}{0.3pt}\par
% loc: opening and Hom insertion
$K\subset\mathbb C$ sous-corps, $M$ motif sur $K$, d'où $T_{\mathrm{DR}}(M)$. Considérons $T_{\mathrm{DR}}(M\unsure{\leftrightharpoons|4|{}^{G}})$\inserted{$=\operatorname{Hom}(1,M)\otimes_{\mathbb Q}K$}, on voit que c'est un sous-module de $T_{\mathrm{DR}}(M)$, \deleted{qui est}\inserted{et que}
% loc: clauses a and b
\par a) contenu dans $T_{\mathrm{DR}}(M)_{\mathbb C}^{\unsure{00|\circ\circ}}$.
\par b) \inserted{$\operatorname{Hom}(1,M)\otimes_{\mathbb Q}(K\cap\mathbb R)$ est} contenu dans $\bigl(T_{\mathrm{DR}}(M)_{\mathbb C}\bigr)_{\mathbb R}$ \deleted{[et où $\operatorname{Hom}(1,M)\otimes_{\mathbb Q}(K\cap\mathbb R)$ est contenu dans $\bigl(T_{\mathrm{DR}}(M)_{\mathbb C}\bigr)_{\mathbb R}$]}

% loc: question
A-t-on une réciproque ?

% loc: item 1, opening and group insertions
$1^\circ$) Pas pour \deleted{\unsure{a) tout seul, même si}}\inserted{a) et b) simultanément} $K=\mathbb Q$, si le groupe de galois motivique $G$ \inserted{pour la cat. de motifs $\unsure{\mathcal M|M}$ où on travaille} n'est pas connexe. En effet, si $G$ est \deleted{\unsure{discret}} fini [Cas des motifs effectifs de poids 0] la condition a) est vide, car $T_{\mathrm{DR}}(M)_{\mathbb C}^{00}=T_{\mathrm{DR}}(M)_{\mathbb C}$, \inserted{donc réc. de a) vraie ssi $G=(e)$.} d'autre part \deleted{b) est} la réciproque à b) est vraie ssi l'extension galoisienne $L$ de groupe $G$ de $\mathbb Q$ définie par $\mathcal M$ est \uline{tot. réelle} [on teste avec le $G$-module $M=\mathbb Q[\unsure{C|G}]$, qui donne $T_{\mathrm{DR}}(M)=L$ \ldots] ; donc dans ce cas on a une réciproque à a) et b).

% loc: connected group and embeddings
Il faut donc se borner à un $G$ connexe, ce qui suggère de regarder les réciproques pour le cas $K$ alg. clos, \inserted{p.ex. $K=\overline{\mathbb Q}$} On peut de plus renforcer la condition, en regardant \uline{toutes} les immersions $K\hookrightarrow\mathbb C$.

% loc: item 2, introductory lines
$2^\circ$) La réciproque à a) est fausse, \deleted{\unsure{même si}} ($=$ pour $\overline{\mathbb Q}$). il suffit de prendre un groupe de Galois motivique $G$ qui est un \uline{tore} de dim $\geq 3$. On voit qu'on peut trouver \inserted{\unsure{sur|hom}}
% loc: pair of arrows, target, and marginal additions
$j_1,j_2:{\unsure{\mathbb G_{\mathrm m\,K}|\mathrm G_{\mathrm m\,K}}}^{\unsure{p|\prime}}\rightrightarrows G'_{\deleted{\unsure{K}}\unsure{K|\mathbb C}}$\inserted{$(=\operatorname{ad}(P)$, car $G$ \unsure{commut.})} qui défnt \inserted{(pour tt $K\longrightarrow\mathbb C$)} la bigraduation de
% loc: struck margin beside the two arrows
\margin{\deleted{\uline{NB} $j_1$ et $j_2$ \unsure{proviennent} d'\unsure{un choix} de \unsure{deux} \unsure{immersions} \unsure{dans} $\mathbb C$, et \illegible{1}}}
% loc: scalar extensions and torsor
$\deleted{\unsure{G}}V_{\mathbb C}\simeq V'\otimes_K\mathbb C$ \quad [$V'=P\mathbin{\overset{G_K}{\deleted{\unsure{\times}}}}V$, $P\unsure{\simeq|=}\underline{\operatorname{Isom}}_{\otimes}(T_{\mathrm B}\otimes K,T_{\mathrm{DR}})$]
% loc: degree-zero representation
Donc une représentation \inserted{fidèle} de degré 0 de $G$ \unsure{donne|donnée}, $\deleted{\unsure{V^{G'}}}$ si $T$ est l'image de $(j_1,j_2):\mathbb G_{\mathrm m\,K}\times\mathbb G_{\mathrm m\,K}\unsure{-|\longrightarrow}G'$, un $V'$ pour lequel
% loc: final equality and inequality
${V'}^{T}\inserted{(=T_{\mathrm{DR}}\cap(T_{\mathrm{DR}}\deleted{\unsure{)}}{}_{\mathbb C}^{00})}$
\inserted{\deleted{\unsure{quelque}}} est $\ne {V'}^{G'}$.
% loc: paper watermark, read sideways and compared across sheets
% printed form
% paper watermark, visible letters in the direction of the paper lettering
\par\noindent\textsc{Extra Strong}\\
\unsure{Maunoury|Maunouy} 1808
% archive sheet number: 50

